\documentclass[12pt,reqno]{article}

\usepackage[usenames]{color}
\usepackage{amssymb,amsmath,amsthm,amsfonts}
\usepackage[colorlinks=true,linkcolor=webgreen,filecolor=webbrown,citecolor=webgreen,urlcolor=webgreen]{hyperref}
\definecolor{webgreen}{rgb}{0,.5,0}
\definecolor{webbrown}{rgb}{.6,0,0}

\setlength{\textwidth}{6.5in}
\setlength{\oddsidemargin}{.1in}
\setlength{\evensidemargin}{.1in}
\setlength{\topmargin}{-.1in}
\setlength{\textheight}{8.4in}

\newcommand{\seqnum}[1]{\href{https://oeis.org/#1}{\rm \underline{#1}}}
\newcommand{\lean}[1]{\texttt{#1}}

\theoremstyle{plain}
\newtheorem{theorem}{Theorem}
\newtheorem{corollary}[theorem]{Corollary}
\newtheorem{lemma}[theorem]{Lemma}
\newtheorem{proposition}[theorem]{Proposition}
\newtheorem{conjecture}{Conjecture}
\theoremstyle{definition}
\newtheorem*{definition}{Definition}
\theoremstyle{remark}
\newtheorem{remark}[theorem]{Remark}

\begin{document}

\begin{center}
\vskip 1cm{\LARGE\bf
A Conjecture on the Least Deficient\\
\vskip .1in
Multiplier Is Equivalent to the\\
\vskip .1in
Almost Perfect Number Conjecture}
\vskip 1cm
\large
Hao Li\\
Metro Vancouver, BC\\
Canada\\
\href{mailto:haoli.leo@outlook.com}{\tt haoli.leo@outlook.com}
\end{center}

\vskip .2 in

\begin{abstract}
For $n \ge 2$ let $a(n)$ be the smallest deficient number $k$ such that $kn$ is
not deficient; this is sequence \seqnum{A215926} in the OEIS. In 2012 Marcus
conjectured that $a(n)$ is always $1$, $3$, or a power of $2$, and that for
$m \ge 1$ the value $2^m$ first occurs at the least prime exceeding $2^m$. We
prove that the first conjecture is \emph{equivalent} to the classical conjecture
that every almost perfect number (a number $k$ with $\sigma(k) = 2k-1$) is a
power of $2$, and that the second conjecture follows from it. Unconditionally,
$a(n) \in \{1,2,3\}$ for even $n$, and for odd $n$ the value $a(n)$ is at most
$n$ and is either a power of $2$ or an almost perfect number that is not a
power of $2$. Consequently both conjectures hold in the range covered by any
search for almost perfect numbers; a computer search up to $10^{16}$ gives the
first conjecture for all $n \le 10^{16}$ and the second for all $m \le 53$. All
theorems have been formally verified in Lean~4 with Mathlib; the computer search
has not. The proofs and the formalization were produced with an AI system,
as described in the paper.
\end{abstract}

\section{Introduction}

Let $\sigma(n)$ denote the sum of the positive divisors of $n$, and write
$h(n) = \sigma(n)/n$ for the abundancy index. A positive integer $n$ is
\emph{deficient} if $\sigma(n) < 2n$, \emph{perfect} if $\sigma(n) = 2n$, and
\emph{abundant} if $\sigma(n) > 2n$; we call $n$ \emph{non-deficient} if it is
perfect or abundant. The deficient numbers form sequence \seqnum{A005100} and
the non-deficient numbers form sequence \seqnum{A023196} in the On-Line
Encyclopedia of Integer Sequences (OEIS)~\cite{oeis}.

\begin{definition}
For an integer $n \ge 2$ let
\[
  T(n) = \{\, k \ge 1 : k \text{ is deficient and } kn \text{ is non-deficient} \,\},
  \qquad a(n) = \min T(n).
\]
\end{definition}

We show below that $T(n)$ is nonempty for every $n \ge 2$, so $a(n)$ is well
defined. The sequence $(a(n))_{n \ge 2}$ is \seqnum{A215926}, contributed to the
OEIS by Michel Marcus in August 2012. It begins
\[
 3, 2, 3, 4, 1, 4, 3, 2, 2, 8, 1, 8, 2, 2, 3, 16, 1, 16, 1, 2, 3, 16, 1, 4, 3, 2, 1, 16, 1, 16, \ldots
\]
The entry records two conjectures. Let $p_m$ denote the least prime greater than
$2^m$ (sequence \seqnum{A014210}).

\begin{conjecture}[Marcus, \seqnum{A215926}]\label{conj:one}
For every $n \ge 2$, the number $a(n)$ is $1$, $3$, or a power of $2$.
\end{conjecture}

\begin{conjecture}[Marcus, \seqnum{A215926}]\label{conj:two}
For every $m \ge 1$, the least $n \ge 2$ with $a(n) = 2^m$ is $n = p_m$.
\end{conjecture}

Both statements appear as open research problems in the
\emph{formal-conjectures} repository~\cite{fc}, a collection of formalized
statements of open problems, and we use the formal statements given there.

A positive integer $k$ is \emph{almost perfect} if $\sigma(k) = 2k-1$. Every
power of $2$ is almost perfect, since $\sigma(2^j) = 2^{j+1}-1$. Whether there
are others is a well-known open problem \cite[Section B2]{guy}.

\begin{conjecture}[Almost perfect number conjecture]\label{conj:ap}
Every almost perfect number is a power of $2$.
\end{conjecture}

We call an almost perfect number that is not a power of $2$ \emph{exotic}.
Kishore~\cite[Proposition~4]{kishore78} showed that an odd exotic almost perfect
number has at least six distinct prime factors, and Antalan and
Dris~\cite{antalan-dris} obtained restrictions on even ones; see also
\cite{dris}. Our results are the following.

\begin{theorem}\label{thm:structure}
Let $n \ge 2$.
\begin{itemize}
\item[(a)] If $n$ is even, then $a(n) \in \{1,2,3\}$.
\item[(b)] If $n$ is odd, then $a(n) \le 2^{\lfloor \log_2 n \rfloor} \le n$, and
$a(n)$ is either a power of $2$ or an exotic almost perfect number.
\end{itemize}
\end{theorem}

\begin{theorem}\label{thm:equiv}
Conjecture~\ref{conj:one} holds if and only if Conjecture~\ref{conj:ap} holds.
More precisely, if $k$ is the least exotic almost perfect number, then there is
an odd squarefree integer $n$ with $a(n) = k$.
\end{theorem}

\begin{theorem}\label{thm:two}
Let $m \ge 1$.
\begin{itemize}
\item[(a)] If $n \ge 2$ and $a(n) = 2^m$, then $n \ge p_m$.
\item[(b)] If there is no exotic almost perfect number less than $2^m$, then
$a(p_m) = 2^m$.
\end{itemize}
In particular, Conjecture~\ref{conj:ap} implies Conjecture~\ref{conj:two}.
\end{theorem}

Thus Conjecture~\ref{conj:one} is exactly as hard as a classical open problem,
while a counterexample to either conjecture of Marcus would have to come from an
exotic almost perfect number. Theorems~\ref{thm:structure} and~\ref{thm:two}
also turn any numerical bound on exotic almost perfect numbers into a range of
validity. The arguments are short and elementary; the point of this note is the
link between the two problems.

\begin{corollary}\label{cor:bound}
Suppose that there is no exotic almost perfect number $k \le N$. Then
Conjecture~\ref{conj:one} holds for all $2 \le n \le N$, and
Conjecture~\ref{conj:two} holds for all $m \ge 1$ with $2^m \le N+1$.
\end{corollary}

A direct search by a C program (Section~\ref{sec:computation}), which is
outside the formal verification, finds no exotic almost perfect number
$k \le 10^{16}$. Hence Conjecture~\ref{conj:one} holds for all
$n \le 10^{16}$ and Conjecture~\ref{conj:two} holds for all $m \le 53$.

All the theorems above have been formalized and checked in the Lean~4 proof
assistant~\cite{lean4} using the Mathlib library~\cite{mathlib}; the numerical
search has not. Details are given in Section~\ref{sec:lean}.

\section{Preliminaries}

We use the following standard facts.

\begin{lemma}\label{lem:sigma}
Let $k, n, t$ be positive integers.
\begin{itemize}
\item[(a)] $\sigma(kn) \le \sigma(k)\sigma(n)$, with equality if $\gcd(k,n)=1$.
\item[(b)] $\sigma(nt) \ge t\,\sigma(n)$; equivalently $h(nt) \ge h(n)$.
\item[(c)] $\sigma(2^j) = 2^{j+1}-1$ for $j \ge 0$, and $\sigma(p) = p+1$ for a prime $p$.
\item[(d)] If $r$ is a divisor of $n$ with $r < n$, then $\sigma(n) \ge n + r$.
\end{itemize}
\end{lemma}

\begin{proof}
(a) Every divisor of $kn$ can be written as $d_1 d_2$ with $d_1 \mid k$ and
$d_2 \mid n$, so the divisors of $kn$ are among the terms of the expanded
product $\sigma(k)\sigma(n)$; multiplicativity of $\sigma$ gives the case of
equality. (b) The numbers $dt$ with $d \mid n$ are distinct divisors of $nt$.
(c) and (d) are immediate.
\end{proof}

Since $\sigma(2^j) = 2^{j+1}-1 < 2^{j+1}$, every power of $2$ is deficient. For
odd $n$, Lemma~\ref{lem:sigma}(a) and (c) give
\begin{equation}\label{eq:pow2}
  2^j \in T(n) \iff (2^{j+1}-1)\,\sigma(n) \ge 2^{j+1} n
  \iff h(n) \ge \frac{2^{j+1}}{2^{j+1}-1}.
\end{equation}

\section{The structure of \texorpdfstring{$a(n)$}{a(n)}}

\begin{proof}[Proof of Theorem~\ref{thm:structure}(a)]
Let $n$ be even. Then $6 \mid 3n$, so $h(3n) \ge h(6) = 2$ by
Lemma~\ref{lem:sigma}(b), and $3n$ is non-deficient. As $3$ is deficient,
$3 \in T(n)$ and $a(n) \le 3$.
\end{proof}

The main step is the following lemma.

\begin{lemma}\label{lem:key}
Let $n \ge 3$ be odd and let $k = a(n)$. If $k$ is not a power of $2$, then
$\sigma(k) = 2k-1$.
\end{lemma}

\begin{proof}
Choose $j \ge 0$ with $2^j < k < 2^{j+1}$. Since $2^j < a(n)$ and $2^j$ is
deficient, $2^j \notin T(n)$, so by~\eqref{eq:pow2}
\begin{equation}\label{eq:upper}
  h(n) < \frac{2^{j+1}}{2^{j+1}-1}.
\end{equation}
On the other hand $kn$ is non-deficient, so by Lemma~\ref{lem:sigma}(a)
\begin{equation}\label{eq:lower}
  2kn \le \sigma(kn) \le \sigma(k)\,\sigma(n).
\end{equation}
Suppose that $\sigma(k) \le 2k-2$. Then \eqref{eq:lower} gives
$h(n) \ge k/(k-1)$, and with \eqref{eq:upper} we get
$k/(k-1) < 2^{j+1}/(2^{j+1}-1)$. Since $x \mapsto x/(x-1)$ is decreasing for
$x > 1$, this forces $k > 2^{j+1}$, a contradiction. Hence $\sigma(k) \ge 2k-1$,
and since $k$ is deficient, $\sigma(k) = 2k-1$.
\end{proof}

\begin{proof}[Proof of Theorem~\ref{thm:structure}(b)]
Let $n \ge 3$ be odd and put $J = \lfloor \log_2 n \rfloor$, so that
$n + 1 \le 2^{J+1}$. By Lemma~\ref{lem:sigma}(d), $\sigma(n) \ge n+1$, hence
\[
 (2^{J+1}-1)\,\sigma(n) \ge (2^{J+1}-1)(n+1) = 2^{J+1}n + 2^{J+1} - (n+1) \ge 2^{J+1} n ,
\]
and $2^J \in T(n)$ by~\eqref{eq:pow2}. Thus $T(n) \neq \emptyset$ and
$a(n) \le 2^J \le n$. The remaining assertion is Lemma~\ref{lem:key}.
\end{proof}

Theorem~\ref{thm:structure} shows that Conjecture~\ref{conj:ap} implies
Conjecture~\ref{conj:one}, and it proves the first half of
Corollary~\ref{cor:bound}: if $2 \le n \le N$ and no exotic almost perfect number
is at most $N$, then $a(n) \in \{1,2,3\}$ or $a(n)$ is a power of $2$.

\begin{remark}\label{rem:formula}
The proof gives an explicit description of $a(n)$. If $n$ is even, then
$a(n) = 1$ if $n$ is non-deficient, $a(n) = 2$ if $n$ is deficient and $2n$ is
not, and $a(n) = 3$ otherwise. If $n \ge 3$ is odd and no exotic almost perfect
number is at most $n$, then by~\eqref{eq:pow2} we have $a(n) = 2^j$, where
$j \ge 0$ is least with $2^{j+1} \ge \sigma(n)/s(n)$ and
$s(n) = \sigma(n) - n$. For example $a(11) = 8$ because
$\sigma(11)/s(11) = 12$.
\end{remark}

\section{The converse}

\begin{proof}[Proof of Theorem~\ref{thm:equiv}]
One direction was just proved. For the other, suppose that exotic almost perfect
numbers exist and let $k$ be the least one. Put $c = \sigma(k) = 2k-1$, choose
$j$ with $2^j < k < 2^{j+1}$, and put $B = 2^{j+1}-1 = \sigma(2^j)$. Then
\begin{equation}\label{eq:kBc}
  k \le B \qquad\text{and}\qquad B + 2 \le c ,
\end{equation}
the second because $k \ge 2^j+1$. Note that $k \ge 5$, since $1, 2, 4$ are
powers of $2$ and $3$ is not almost perfect.

Let $q_1 < q_2 < \cdots$ be the primes that are at least $(c+1)^2$, and let
$n_t = q_1 q_2 \cdots q_t$, with $n_0 = 1$. By Lemma~\ref{lem:sigma},
\[
  h(n_t) = \prod_{i=1}^{t} \Bigl(1 + \frac{1}{q_i}\Bigr) \ge 1 + \sum_{i=1}^{t} \frac{1}{q_i},
\]
which is unbounded because the sum of the reciprocals of the primes
diverges~\cite[Theorem 19]{hw}. Since $h(n_0) = 1 < (c+1)/c$, there is a least
$t \ge 1$ with $h(n_t) \ge (c+1)/c$. Put $n = n_t$ and $q = q_t$. We claim that
$a(n) = k$.

\emph{Step 1: $k \in T(n)$.} The number $k$ is deficient. All prime factors of
$n$ exceed $k$, so $\gcd(k,n) = 1$ and
\[
  h(kn) = h(k)\,h(n) \ge \frac{c}{k}\cdot\frac{c+1}{c} = \frac{2k}{k} = 2 .
\]

\emph{Step 2: $h(n) < (B+1)/B$.} By the minimality of $t$ we have
$h(n_{t-1}) < (c+1)/c$, so
\[
  h(n) = h(n_{t-1})\Bigl(1 + \frac1q\Bigr) < \frac{c+1}{c}\cdot\frac{q+1}{q} .
\]
It suffices to show that $B(c+1)(q+1) \le c(B+1)q$, that is,
$B(c+1) \le q\,(c - B)$. This holds because $c - B \ge 2$ by~\eqref{eq:kBc} and
$q \ge (c+1)^2 > B(c+1)$.

\emph{Step 3: no deficient $k' < k$ lies in $T(n)$.} Let $k' < k$ be deficient.
Again $\gcd(k',n) = 1$, so by Step~2
\[
  h(k'n) = h(k')\,h(n) < h(k')\cdot\frac{B+1}{B}.
\]
If $\sigma(k') = 2k'-1$, then by the minimality of $k$ we have $k' = 2^i$ with
$i \le j$, and
\[
  h(k')\cdot\frac{B+1}{B} = 2\cdot\frac{2^{i+1}-1}{2^{i+1}}\cdot\frac{2^{j+1}}{2^{j+1}-1} \le 2 .
\]
Otherwise $\sigma(k') \le 2k'-2$, so $h(k') \le 2(k'-1)/k'$, and
\[
  h(k')\cdot\frac{B+1}{B} \le 2\cdot\frac{(k'-1)(B+1)}{k'B} \le 2,
\]
because $(k'-1)(B+1) \le k'B$ is equivalent to $k'-1 \le B$, which holds
by~\eqref{eq:kBc}. In both cases $h(k'n) < 2$, so $k' \notin T(n)$.

Steps 1 and 3 give $a(n) = k$. Since $k$ is not $1$, not $3$, and not a power of
$2$, Conjecture~\ref{conj:one} fails at $n$. By construction $n$ is odd and
squarefree.
\end{proof}

\begin{remark}
The integer $n$ produced by the proof is very large. If there happens to be a
prime $q$ with $2^{j+1} \le q \le 2k-1$, the same argument shows that already
$a(q) = k$.
\end{remark}

\section{The first occurrence of \texorpdfstring{$2^m$}{2\^{}m}}

\begin{proof}[Proof of Theorem~\ref{thm:two}(a)]
Let $m \ge 1$ and $a(n) = 2^m$. If $n$ is even, then $a(n) \le 3$ by
Theorem~\ref{thm:structure}(a), so $m = 1$ and $p_1 = 3$. Since $2$ and $4$ are
deficient, $1, 2 \notin T(2)$ and $a(2) = 3 \neq 2$; hence $n \ge 4 > p_1$.

Now let $n$ be odd. Since $2^{m-1} < a(n)$, we have $2^{m-1} \notin T(n)$, so
$(2^m - 1)\,\sigma(n) < 2^m n$ by~\eqref{eq:pow2}. Let $q$ be the least prime
factor of $n$. By Lemma~\ref{lem:sigma}(d), $\sigma(n) \ge n + n/q$, hence
$(2^m-1)(1 + 1/q) < 2^m$, that is, $q > 2^m - 1$. As $q$ is odd and $2^m$ is
even, $q > 2^m$. Therefore $p_m \le q \le n$.
\end{proof}

\begin{proof}[Proof of Theorem~\ref{thm:two}(b)]
Let $p = p_m$. By Bertrand's postulate \cite[Theorem 418]{hw} there is a prime
in the interval $(2^m, 2^{m+1}]$, and $2^{m+1}$ is not prime, so
$2^m < p < 2^{m+1}$. In particular $p$ is odd and $\sigma(p) = p+1$, so
by~\eqref{eq:pow2}
\begin{equation}\label{eq:prime}
  2^i \in T(p) \iff (2^{i+1}-1)(p+1) \ge 2^{i+1}p \iff 2^{i+1} \ge p+1 .
\end{equation}
Taking $i = m$ shows $2^m \in T(p)$, so $a(p) \le 2^m$. By
Theorem~\ref{thm:structure}(b), $a(p)$ is a power of $2$ or an exotic almost
perfect number. The latter would be an exotic almost perfect number less than
$2^m$, which is excluded by hypothesis. So $a(p) = 2^i$ for some $i \le m$, and \eqref{eq:prime} gives
$2^{i+1} \ge p + 1 > 2^m$, whence $i = m$.
\end{proof}

Parts (a) and (b) together give Conjecture~\ref{conj:two} for every $m$ such
that no exotic almost perfect number is less than $2^m$, which completes the
proof of Corollary~\ref{cor:bound}.

\begin{remark}
We do not know whether Conjecture~\ref{conj:two} implies
Conjecture~\ref{conj:ap}. The proof shows that Conjecture~\ref{conj:two} fails
for a given $m$ exactly when there is an exotic almost perfect number $k$ with
$(p_m+1)/2 \le k < 2^m$: such a $k$ is coprime to $p_m$ and satisfies
$h(k p_m) = (2k-1)(p_m+1)/(k p_m) \ge 2$, so $a(p_m) \le k < 2^m$; conversely, an
exotic value $a(p_m) = k$ must satisfy this inequality.
\end{remark}

\section{A search for almost perfect numbers}\label{sec:computation}

If $\sigma(k) = 2k-1$ then $\sigma(k)$ is odd, so $k$ is a square or twice a
square; that is, $k = 2^a b^2$ with $a \ge 0$ and $b$ odd. For such $k$ the
condition $\sigma(k) = 2k-1$ reads
\[
  (2^{a+1}-1)\,\sigma(b^2) = 2^{a+1} b^2 - 1 .
\]
We tested this identity for every odd $b$ with $3 \le b \le 10^8$ and every
$a \ge 0$ with $2^a b^2 \le 10^{16}$, computing $\sigma(b^2)$ from the
factorization of $b$ obtained with a sieve. There are $170{,}710{,}624$ such
pairs, and none satisfies the identity. Hence there is no exotic almost perfect
number $k \le 10^{16}$. By Corollary~\ref{cor:bound}, Conjecture~\ref{conj:one}
holds for all $n \le 10^{16}$ and Conjecture~\ref{conj:two} holds for all
$m \le 53$, since $2^{53} < 10^{16}$.

This computation is elementary and we make no claim that the bound $10^{16}$ is
new or best possible. It is the only part of this paper that has not been
formally verified; the program is available with the formal proof.

\section{Formal verification}\label{sec:lean}

The results of this paper were formalized in Lean~4 (version 4.36.0) with
Mathlib (version v4.36.0). The formalization is a single file,
\lean{A215926.lean}, of about 800 lines, available in the directory
\lean{a215926} of the repository
\begin{center}
\url{https://github.com/Leo309/lean-open-problems}
\end{center}
The definitions of $a(n)$, of $p_m$, and of almost perfect numbers, as well as
the statements of Conjectures~\ref{conj:one}, \ref{conj:two}, and of the open
question on exotic almost perfect numbers, are copied verbatim from two files
of~\cite{fc}, namely
\begin{center}
\lean{FormalConjectures/OEIS/215926.lean},\\
\lean{FormalConjectures/Wikipedia/AlmostPerfectNumbers.lean}.
\end{center}
The
file contains no \lean{sorry}, and every theorem depends only on the three
standard axioms \lean{propext}, \lean{Classical.choice}, and \lean{Quot.sound}.
The correspondence with the present paper is as follows.

\begin{center}
\begin{tabular}{ll}
Result & Lean theorem \\ \hline
Theorem~\ref{thm:structure}(a) & \lean{a\_even} \\
Lemma~\ref{lem:key}, Theorem~\ref{thm:structure}(b) & \lean{a\_odd}, \lean{two\_pow\_log\_mem}, \lean{a\_le\_self} \\
Theorem~\ref{thm:equiv}, ``if'' & \lean{conjecture\_1\_of\_almostPerfect} \\
Theorem~\ref{thm:equiv}, ``only if'' & \lean{exists\_a\_eq\_of\_exotic}, \lean{almostPerfect\_of\_conjecture\_1} \\
Theorem~\ref{thm:equiv} & \lean{conjecture\_1\_iff\_almostPerfect} \\
Theorem~\ref{thm:two}(a) & \lean{nextPrime\_le\_of\_a\_eq} \\
Theorem~\ref{thm:two}(b) & \lean{a\_nextPrime}, \lean{conjecture\_2\_of\_almostPerfect} \\
Corollary~\ref{cor:bound} & \lean{conjecture\_1\_of\_bound}, \lean{conjecture\_2\_of\_bound} \\
\end{tabular}
\end{center}

The divergence of the sum of the reciprocals of the primes and Bertrand's
postulate enter through the Mathlib lemmas
\begin{center}
\lean{one\_half\_le\_sum\_primes\_ge\_one\_div},\\
\lean{Nat.exists\_prime\_lt\_and\_le\_two\_mul}.
\end{center}
In addition,
the file proves Conjecture~\ref{conj:one} for $2 \le n < 256$ and
Conjecture~\ref{conj:two} for $1 \le m \le 8$ unconditionally, by letting the
Lean kernel check that there is no exotic almost perfect number below $256$
(\lean{conjecture\_1\_small}, \lean{conjecture\_2\_small}). This is only a
consistency check of the formal definitions and gives no new information, since
the table of $a(n)$ for $n \le 10000$ in the OEIS entry already covers a larger
range.
Remark~\ref{rem:formula}, the two other remarks, the statement that $n$ is odd
and squarefree in Theorem~\ref{thm:equiv}, and the search of
Section~\ref{sec:computation} are not part of the formalization.

\section{Use of artificial intelligence}

The mathematical results of this paper, their proofs, the Lean formalization,
the search program, and the first draft of the text were produced by Claude, an
AI system developed by Anthropic, working at the author's request. The link
between sequence \seqnum{A215926} and almost perfect numbers was found by the AI
system while it searched for counterexamples to Conjecture~\ref{conj:one}. The
correctness of the theorems does not rest on the AI system: it rests on the
proofs given here and on their verification by the Lean kernel. The author has
checked that the formal statements agree with the definitions in the OEIS and
in~\cite{fc}, has read the proofs, and takes full responsibility for the
contents of the paper.

\section{Acknowledgments}

The author thanks the contributors and editors of the OEIS, and the maintainers
of Mathlib and of the formal-conjectures repository.

\begin{thebibliography}{9}

\bibitem{antalan-dris}
J. R. M. Antalan and J. A. B. Dris, Some new results on even almost perfect
numbers which are not powers of two, preprint, 2016. Available at
\url{https://arxiv.org/abs/1602.04248}.

\bibitem{dris}
J. A. B. Dris, A criterion for deficient numbers using the abundancy index and
deficiency functions, preprint, 2013. Available at
\url{https://arxiv.org/abs/1308.6767}.

\bibitem{fc}
Google DeepMind, formal-conjectures, a collection of formalized statements of
conjectures in Lean, 2026. Available at
\url{https://github.com/google-deepmind/formal-conjectures}.

\bibitem{guy}
R. K. Guy, \emph{Unsolved Problems in Number Theory}, 3rd edition, Springer,
2004.

\bibitem{hw}
G. H. Hardy and E. M. Wright, \emph{An Introduction to the Theory of Numbers},
4th edition, Clarendon Press, Oxford, 1960.

\bibitem{kishore78}
M. Kishore, Odd integers $N$ with five distinct prime factors for which
$2 - 10^{-12} < \sigma(N)/N < 2 + 10^{-12}$, \emph{Math. Comp.} \textbf{32}
(1978), 303--309.

\bibitem{lean4}
L. de Moura and S. Ullrich, The Lean 4 theorem prover and programming language,
in \emph{Automated Deduction -- CADE 28}, Lect. Notes in Comput. Sci.,
Vol.~12699, Springer, 2021, pp.~625--635.

\bibitem{mathlib}
The mathlib Community, The Lean mathematical library, in \emph{Proceedings of
the 9th ACM SIGPLAN International Conference on Certified Programs and Proofs
(CPP 2020)}, ACM, 2020, pp.~367--381.

\bibitem{oeis}
OEIS Foundation Inc., The On-Line Encyclopedia of Integer Sequences, 2026.
Available at \url{https://oeis.org}.

\end{thebibliography}

\bigskip
\hrule
\bigskip

\noindent 2020 {\it Mathematics Subject Classification}:
Primary 11A25; Secondary 11B83, 68V20.

\noindent \emph{Keywords:}
deficient number, abundant number, almost perfect number, sum of divisors,
abundancy index, formal proof, Lean.

\bigskip
\hrule
\bigskip

\noindent (Concerned with sequences
\seqnum{A005100},
\seqnum{A014210},
\seqnum{A023196}, and
\seqnum{A215926}.)

\end{document}
